\documentclass[12pt]{article}
\usepackage[utf8]{inputenc}
\usepackage[spanish]{babel}
\usepackage{amsmath}
\usepackage{amssymb}
\usepackage{geometry}
\geometry{a4paper, margin=2cm}

\begin{document}

\begin{center}
    \textbf{SOLUCIONARIO: EXAMEN DE RECUPERATORIO}
    \\\textbf{Transferencia de Calor - Escuela 4117}
\end{center}

\section*{Ejercicio 1: Conversiones}

\begin{itemize}
    \item[a)] $F = \frac{9}{5}(10) + 32 \rightarrow 18 + 32 = \mathbf{50^\circ F}$
    \item[b)] $C = \frac{5}{9}(68 - 32) \rightarrow \frac{5}{9}(36) \rightarrow 5 \cdot 4 = \mathbf{20^\circ C}$
    \item[c)] $K = 27 + 273 = \mathbf{300\text{ K}}$
    \item[d)] $10\text{ BTU} \cdot 252\text{ cal/BTU} = \mathbf{2520\text{ cal}}$
\end{itemize}

\section*{Ejercicio 2: Calorimetría}

\begin{itemize}
    \item \textbf{Fórmula:} $Q = m \cdot C_e \cdot \Delta T$
    \item \textbf{Cálculo:} $Q = 5\text{ kg} \cdot 1\text{ kcal/kg}^\circ\text{C} \cdot 20^\circ\text{C} = \mathbf{100\text{ kcal}}$
\end{itemize}

\section*{Ejercicio 3: Conducción}

\begin{itemize}
    \item \textbf{Fórmula:} $Q = \frac{k \cdot A \cdot \Delta T}{L}$
    \item \textbf{Cálculo:} $Q = \frac{1.5 \cdot 2 \cdot 30}{0.1} = \frac{90}{0.1} = \mathbf{900\text{ kcal/h}}$
\end{itemize}

\vspace{1cm}
\noindent \textit{Nota Pedagógica: Si el alumno llega a un resultado muy diferente a estos números redondos, es la señal clara para que revisen el despeje o la conversión de unidades (especialmente los $10\text{ cm} \rightarrow 0.1\text{ m}$).}

\end{document}